\SourcePage{619}{622}
\section{Born's formula}

The scattering cross-section can be calculated in general form in the very
important case where the scattering field may be treated as a
perturbation\footnote{In the general theory developed in \S~123, this
approximation corresponds to all the phases $\delta_l$ being small; in
addition, these phases must be calculable from the Schrödinger equation with
the potential energy treated as a perturbation (see Problem 4).}. Section 45
showed that this is possible if at least one of the two conditions
\begin{equation}
 |U|\ll\frac{\hbar^2}{ma^2}
 \tag{126.1}
\end{equation}
or
\begin{equation}
 |U|\ll\frac{\hbar v}{a}=\frac{\hbar^2ka}{ma^2}
 \tag{126.2}
\end{equation}
is satisfied. Here $a$ is the range of the field $U(r)$, and $U$ denotes its
order of magnitude in the principal region where it is present. When the first

\SourcePage{620}{623}
condition holds, the approximation applies at every velocity. The second
condition shows that it applies in any event to sufficiently fast particles.

Following \S~45, write the wave function as
$\psi=\psi^{(0)}+\psi^{(1)}$, where $\psi^{(0)}=e^{i\mathbf{kr}}$ describes the
incident particle with wave vector $\mathbf k=\mathbf p/\hbar$. Formula (45.3)
gives
\begin{equation}
 \psi^{(1)}(x,y,z)=-\frac{m}{2\pi\hbar^2}\int
 \frac{U(x',y',z')e^{i(\mathbf{kr'}+kR)}}{R}\,dV'.
 \tag{126.3}
\end{equation}
Choose the scattering center as origin, introduce the radius vector
$\mathbf R_0$ to the observation point of $\psi^{(1)}$, and let $\mathbf n'$
be the unit vector along $\mathbf R_0$. If the radius vector of the volume
element $dV'$ is $\mathbf r'$, then $\mathbf R=\mathbf R_0-\mathbf r'$. Far
from the center, $R_0\gg r'$, and hence
\[
 R=|\mathbf R_0-\mathbf r'|\approx R_0-\mathbf r'\mathbf n'.
\]
Insertion in (126.3) yields the asymptotic expression
\[
 \psi^{(1)}\approx-\frac{m}{2\pi\hbar^2}\frac{e^{ikR_0}}{R_0}
 \int U(\mathbf r')e^{i(\mathbf k-\mathbf k')\mathbf r'}dV',
\]
where $\mathbf k'=k\mathbf n'$ is the wave vector after scattering. Comparison
with the definition (123.3) gives the scattering amplitude
\begin{equation}
 f=-\frac{m}{2\pi\hbar^2}\int Ue^{-i\mathbf{qr}}dV,
 \tag{126.4}
\end{equation}
after relabeling the integration variables and introducing
\begin{equation}
 \mathbf q=\mathbf k'-\mathbf k,
 \tag{126.5}
\end{equation}
whose magnitude is
\begin{equation}
 q=2k\sin\frac\theta2,
 \tag{126.6}
\end{equation}
with $\theta$ the angle between $\mathbf k$ and $\mathbf k'$, that is, the
scattering angle.

Finally, taking the squared modulus of the amplitude gives the cross-section
into the solid-angle element $do$:
\begin{equation}
 d\sigma=\frac{m^2}{4\pi^2\hbar^4}
 \left|\int Ue^{-i\mathbf{qr}}dV\right|^2do.
 \tag{126.7}
\end{equation}

\SourcePage{621}{624}Thus, scattering that changes the momentum by $\hbar\mathbf q$ is governed by the squared modulus of the matching Fourier component of the field $U$.

Formula~(126.7) was first obtained by \emph{Born} (M.~Born, 1926); this approximation is often called the \emph{Born approximation} in collision theory.


In this approximation the relation
\begin{equation}
 f(\mathbf k,\mathbf k')=f^*(\mathbf k',\mathbf k)
 \tag{126.8}
\end{equation}
holds between the amplitudes of the direct and reverse scattering processes in
the literal sense: the initial and final momenta are interchanged, without
reversal of their signs as under time reversal. Scattering therefore has an
additional symmetry besides the reciprocity theorem (125.12). This symmetry is
closely connected with the small scattering amplitudes of perturbation theory
and follows directly from the unitarity condition (125.8) if its integral term,
quadratic in $f$, is neglected\footnote{It follows that this property already
disappears in the second approximation of perturbation theory. We shall verify
this directly in \S~130 in connection with (130.13).}.

Formula (126.7) can also be obtained in another way, though this leaves the
phase of the scattering amplitude undetermined. Start from the general formula
(43.1) for the transition probability between continuous-spectrum states,
\[
 dw_{fi}=\frac{2\pi}{\hbar}|U_{fi}|^2\delta(E_f-E_i)d\nu_f.
\]
Here it is applied to the transition from an incident free-particle state of
momentum $\mathbf p$ to a state of momentum $\mathbf p'$ scattered into the
solid-angle element $do'$. Choose $d^3p'/(2\pi\hbar)^3$ as the state interval
$d\nu_f$. With $E_f-E_i=(p'^2-p^2)/2m$, we find
\begin{equation}
 dw_{\mathbf p'\mathbf p}=\frac{4\pi m}{\hbar}|U_{\mathbf p'\mathbf p}|^2
 \delta(p'^2-p^2)\frac{d^3p'}{(2\pi\hbar)^3}.
 \tag{126.9}
\end{equation}

The incident and scattered wave functions are plane waves. Since the state
interval was chosen as an element of $\mathbf p'/(2\pi\hbar)$ space, the final
wave function must be normalized to a delta-function of
$\mathbf p'/(2\pi\hbar)$:
\begin{equation}
 \psi_{\mathbf p'}=e^{i\mathbf p'\mathbf r/\hbar}.
 \tag{126.10}
\end{equation}

\SourcePage{622}{625}
Normalize the initial wave function to unit current density:
\begin{equation}
 \psi_{\mathbf p}=\sqrt{\frac mp}\,e^{i\mathbf p\mathbf r/\hbar}.
 \tag{126.11}
\end{equation}
Then (126.9) has the dimension of area and is the differential scattering
cross-section.

The delta-function in (126.9) means that $p'=p$, as required for elastic
scattering. It may be eliminated by using spherical coordinates in momentum
space, replacing $d^3p'$ by
$p'^2dp' do'=(1/2)p'd(p'^2)do'$, and integrating over $d(p'^2)$. This amounts
to replacing $p'$ by $p$ in the integrand and gives
\[
 d\sigma=\frac{4\pi^2}{\hbar^4}p^2
 \left|\int\psi_{\mathbf p'}^*U\psi_{\mathbf p}dV\right|^2do'.
\]
Substitution of (126.10) and (126.11) returns (126.7).

In the form (126.7), the formula applies to a field $U(x,y,z)$ depending on
coordinates in any combination, not merely on $r$. For a central field $U(r)$,
however, it can be transformed further.

In the integral
\[
 \int U(r)e^{-i\mathbf{qr}}dV
\]
use spherical coordinates $r,\vartheta,\varphi$, taking the polar axis along
$\mathbf q$; $\vartheta$ distinguishes the polar angle from the scattering
angle $\theta$. Integration over $\vartheta$ and $\varphi$ gives
\[
 \int_0^\infty\!\int_0^\pi\!\int_0^{2\pi}
 U(r)e^{iqr\cos\vartheta}r^2\sin\vartheta\,d\varphi\,d\vartheta\,dr
 =4\pi\int_0^\infty U(r)\frac{\sin qr}{qr}r^2dr.
\]
Thus the scattering amplitude in a central field is
\begin{equation}
 f=-\frac{2m}{\hbar^2}\int_0^\infty U(r)\frac{\sin qr}{q}r\,dr.
 \tag{126.12}
\end{equation}

\SourcePage{623}{626}
At $\theta=0$, and hence $q=0$, this integral diverges if $U(r)$ decreases at
infinity as $1/r^3$ or more slowly, in agreement with the general results of
\S~124.

Notice an interesting point. The momentum $p$ and scattering angle $\theta$
enter (126.12) only through $q$. In the Born approximation the cross-section
therefore depends on $p$ and $\theta$ only in the combination
$p\sin(\theta/2)$.

Returning to arbitrary fields $U(x,y,z)$, consider the limiting cases of small
($ka\ll1$) and large ($ka\gg1$) velocities.

At small velocities one may put $e^{-i\mathbf{qr}}\approx1$ in (126.4), so
that
\begin{equation}
 f=-\frac{m}{2\pi\hbar^2}\int U\,dV,
 \tag{126.13}
\end{equation}
and, for $U=U(r)$,
\begin{equation}
 f=-\frac{2m}{\hbar^2}\int_0^\infty U(r)r^2dr.
 \tag{126.14}
\end{equation}
The scattering is isotropic and independent of velocity, consistently with
the general results of \S~132.

In the opposite, high-velocity limit the scattering is strongly anisotropic
and directed forward, within a narrow cone of aperture
$\Delta\theta\sim1/ka$. Outside this cone $q$ is large, the factor
$e^{-i\mathbf{qr}}$ oscillates rapidly, and its integral multiplied by the
slowly varying function $U$ is nearly zero.

The decrease of the cross-section at large $q$ is not universal, but depends
on the particular field. If $U(r)$ has a singularity at $r=0$ or at some other
real value of $r$, the neighborhood of that point dominates the integral in
(126.12), and the decrease follows a power law. The same holds when $U(r)$ is
nonsingular but not even: the neighborhood of $r=0$ then gives the main
contribution. If $U(r)$ is even, the integration may formally be extended to
negative $r$, along the entire real axis. Provided $U(r)$ has no singularities
on that axis, the contour can then be shifted into the complex plane until it
is caught by the nearest complex singularity. The integral consequently
decreases exponentially at large $q$.

\SourcePage{624}{627}
It must be remembered, however, that the Born approximation is generally
unsuitable for calculating this exponentially small quantity (see also
\S~131).

Although the differential cross-section inside the cone
$\Delta\theta\sim1/ka$ is essentially independent of velocity, the shrinking
aperture makes the total cross-section decrease at high energies, if
$\int d\sigma$ converges at all. It decreases with the solid angle cut out by
the cone, proportionally to $(\Delta\theta)^2\sim1/k^2a^2$, and hence
inversely with energy.

In many applications of collision theory, scattering is characterized by
\begin{equation}
 \sigma_{\mathrm{tr}}=\int(1-\cos\theta)d\sigma,
 \tag{126.15}
\end{equation}
often called the \emph{transport cross-section}. Arguments like those above
show that at high velocities it is inversely proportional to the square of the
energy.

\subsection*{Problems}

\ProblemHeading{1.} Find in the Born approximation the scattering
cross-section of a spherical potential well: $U=-U_0$ for $r<a$, and $U=0$ for
$r>a$.

\SolutionHeading The integral in (126.12) gives
\[
 d\sigma=4a^2\left(\frac{mU_0a^2}{\hbar^2}\right)^2
 \frac{(\sin qa-qa\cos qa)^2}{(qa)^6}\,do.
\]
Integration over all angles, conveniently performed with
$q=2k\sin(\theta/2)$ and $do=2\pi q\,dq/k^2$, gives
\[
 \sigma=\frac{2\pi}{k^2}\left(\frac{mU_0a^2}{\hbar^2}\right)^2
 \left[1-\frac1{(2ka)^2}+\frac{\sin4ka}{(2ka)^3}
 -\frac{\sin^22ka}{(2ka)^4}\right].
\]
In the limiting cases,
\[
 \sigma=\begin{cases}
 \displaystyle\frac{16\pi a^2}{9}
 \left(\frac{mU_0a^2}{\hbar^2}\right)^2,&ka\ll1,\\[4pt]
 \displaystyle\frac{2\pi}{k^2}
 \left(\frac{mU_0a^2}{\hbar^2}\right)^2,&ka\gg1.
 \end{cases}
\]

\ProblemHeading{2.} The same for the field $U=U_0e^{-r^2/a^2}$.

\SolutionHeading It is convenient to use (126.7), taking $\mathbf q$ along
one coordinate axis. The result is

\SourcePage{625}{628}
\[
 d\sigma=\frac{\pi a^2}{4}\left(\frac{mU_0a^2}{\hbar^2}\right)^2
 e^{-q^2a^2/2}do,
\]
and the total cross-section is
\[
 \sigma=\frac{\pi^2}{2k^2}\left(\frac{mU_0a^2}{\hbar^2}\right)^2
 (1-e^{-2k^2a^2}).
\]
The applicability conditions are (126.1), (126.2), with $U_0$ for $U$.
Moreover, the formula for $d\sigma$ is invalid when the magnitude of the
exponent is large\footnote{This failure of perturbation theory is readily seen
by calculating the second-approximation amplitude (see (130.13) below): its
pre-exponential factor is small compared with the first-approximation
coefficient, but the magnitude of its negative exponent is only half as
large.}.

\ProblemHeading{3.} The same for $U=(\alpha/r)e^{-r/a}$.

\SolutionHeading Evaluation of (126.12) gives
\[
 d\sigma=4a^2\left(\frac{\alpha ma}{\hbar^2}\right)^2
 \frac{do}{(q^2a^2+1)^2}.
\]
The total cross-section is
\[
 \sigma=16\pi a^2\left(\frac{\alpha ma}{\hbar^2}\right)^2
 \frac1{4k^2a^2+1}.
\]
Using $\alpha/a$ for $U$ in (126.1), (126.2), the applicability condition is
$\alpha ma/\hbar^2\ll1$ or $\alpha/\hbar v\ll1$.

\ProblemHeading{4.} Determine the phases $\delta_l$ for scattering in a
central field in the case corresponding to the Born approximation.

\SolutionHeading For the radial wave function $\chi=rR$ in the field $U(r)$
and the free-motion function $\chi^{(0)}$, the equations are (see (32.10))
\[
 \chi''+\left[k^2-\frac{l(l+1)}{r^2}-\frac{2m}{\hbar^2}U\right]\chi=0,
 \qquad
 \chi^{(0)''}+\left[k^2-\frac{l(l+1)}{r^2}\right]\chi^{(0)}=0.
\]
Multiply the first by $\chi^{(0)}$ and the second by $\chi$, subtract, and
integrate over $dr$, using $\chi=0$ at $r=0$. This gives
\[
 \chi'(r)\chi^{(0)}(r)-\chi(r)\chi^{(0)'}(r)
 =\frac{2m}{\hbar^2}\int_0^rU\chi\chi^{(0)}dr.
\]
Treating $U$ as a perturbation, put $\chi\approx\chi^{(0)}$ on the right. At
$r\to\infty$ use (33.12), (33.20) on the left and insert the exact expression
(33.10) in the integral. The result is
\[
 \sin\delta_l\approx\delta_l=-\frac{\pi m}{\hbar^2}
 \int_0^\infty U(r)[J_{l+1/2}(kr)]^2r\,dr.
\]

\SourcePage{626}{629}
The same formula could also be obtained by directly expanding the Born
amplitude (126.4) in Legendre polynomials according to (123.11), for small
$\delta_l$.

\ProblemHeading{5.} In the Born approximation, find the total cross-section
for fast particles ($ka\gg1$) in the field
$U=\alpha/(r^2+a^2)^{n/2}$, where $n>2$.

\SolutionHeading As will emerge, partial amplitudes with large $l$ dominate.
We may therefore use (123.11), replacing the sum over $l$ by an integral. Since
all $\delta_l\ll1$ in the Born approximation,
\begin{equation}
 \sigma\approx\frac{4\pi}{k^2}\int_0^\infty2l\delta_l^2dl.
 \tag{1}
\end{equation}
For large $l$, (124.1) gives
\[
 \delta_l=-\frac{\alpha m}{\hbar^2}\int_{l/k}^\infty
 \frac{dr}{(r^2+a^2)^{n/2}(k^2-l^2/r^2)^{1/2}}.
\]
With $r^2+a^2=(a^2+l^2/k^2)/t$, this reduces to Euler's integral and gives
\begin{equation}
 \delta_l=-\frac{\alpha m}{2\hbar^2k(a^2+l^2/k^2)^{(n-1)/2}}
 \frac{\Gamma(1/2)\Gamma((n-1)/2)}{\Gamma(n/2)}.
 \tag{2}
\end{equation}
The integral (1) is governed by $l\sim ak\gg1$, justifying the assumption.
Evaluation gives
\begin{equation}
 \sigma=\frac{\pi^2}{n-2}
 \left[\frac{\Gamma((n-1)/2)}{\Gamma(n/2)}\right]^2
 \left(\frac{\alpha m}{k\hbar^2a^{n-2}}\right)^2.
 \tag{3}
\end{equation}
By (126.2), the condition for the Born approximation is
$\alpha m/(\hbar^2ka^{n-1})\ll1$. Notice the dependence
$\sigma\sim k^{-2}$, as anticipated by the general statements in the text.

\ProblemHeading{6.} Find the Born-approximation scattering amplitude in two
dimensions, for $U=U(x,z)$ and a particle beam incident along the $z$ axis.

\SolutionHeading Using the note on p.~204 and the familiar asymptotic form of
the Hankel function
\begin{equation}
 H_0^{(1)}(u)\approx\sqrt{\frac2{\pi u}}e^{i(u-\pi/4)}
 \quad\text{as }u\to\infty,
 \tag{4}
\end{equation}
we find, at large distance $\rho_0$ from the field axis (the $y$ axis),
\[
 \psi^{(1)}\sim\frac{f(\varphi)}{\sqrt{\rho_0}}e^{ik\rho_0},
\]
where
\[
 f(\varphi)=-\frac{m}{\hbar^2}\sqrt{\frac{i}{2\pi k}}
 \int U(\boldsymbol\rho)e^{-i\mathbf q\boldsymbol\rho}d^2\rho
\]
is the scattering amplitude. Here $\boldsymbol\rho=(x,z)$ is the
two-dimensional radius vector, $d^2\rho=dx\,dz$, and $\varphi$ is the
scattering angle in the $xz$ plane. In two dimensions the amplitude has the
dimension of the square root of length, while the scattering cross-section
$d\sigma=|f|^2d\varphi$ has the dimension of length.
